\section{总结与延伸}

Reddit 的二十年演化可以看成三条相互缠绕的系统曲线。产品曲线从链接排序扩展到社区、搜索和 AI；治理曲线从“不删除”扩展到分层规则、分级执法和申诉程序；基础设施曲线从单一网站扩展到多云、流量降级、API 与 agent 接口。每一次扩展都在重新解释“由用户驱动”：用户拥有更多表达与工具，同时平台承担更明确的安全、可靠性和可持续责任。

\begin{importantbox}{七条可迁移结论}
\begin{enumerate}
  \item 冷启动的终点是核心循环能在创始人退出后继续，不是注册数达到某个整数。
  \item 网络效应可以掩盖组织与技术债，不能把持续增长当作系统健康证明。
  \item 投票是排序信号，不是全能治理机制；外部性需要更高层规则处理。
  \item 政策文本、执法工具、申诉和透明度共同决定 moderation 质量。
  \item 公共可读、免费 API 与无限商业提取是三种不同访问合同。
  \item 韧性设计要按用户动作排序，并把安全与恢复后核对纳入核心路径。
  \item AI 摘要和 agents 应保留来源、授权和价值回流，否则会消耗社区供给。
\end{enumerate}
\end{importantbox}

\subsection{实践作业：给一个社区平台设计变更方案}

选择一个真实论坛、开源社区或内容平台，提出一项会改变用户合同的变更，例如 API 收费、付费社区、AI 摘要或 agent 写入。作业必须同时覆盖产品、治理、基础设施和迁移程序，不能只给一张商业收益表。读者应能从你的方案中看出受影响群体、风险、回滚条件和成功指标。

\begin{enumerate}
  \item 画出现有核心循环，标记哪些参与者提供内容、审核、分发与收入。
  \item 定义变更后的访问类别、身份、权限、速率和结算规则。
  \item 列出至少三个失败模式，其中一个必须来自激励变化，一个来自基础设施降级。
  \item 设计试点、通知、迁移、申诉和回滚流程，并说明怎样公开复盘。
\end{enumerate}

\begin{knowledgebox}{验收标准}
高质量方案应把“开放”“社区”“公平”“可持续”转成可观察接口和指标。若方案只说“倾听用户”却没有反馈如何进入决策，或只说“收回成本”却没有成本模型和例外规则，说明关键系统仍被口号遮住。
\end{knowledgebox}

\subsection{拓展阅读}

继续学习可沿三条路径展开。第一，对照 Stanford Winter 2025 课程页和本地字幕，建立 Reddit 关键决策的时间线，明确哪些是讲者回忆、哪些需要外部资料核验。第二，阅读一个公开 API 的身份、rate limit、SLA 与弃用政策，评估它是否支持长期生态。第三，选择一个 RAG 或 agent 产品，检查答案是否回链、用户能否撤销授权，以及平台如何处理删除后的缓存和衍生数据。

